\chapter{Plugin Development}

% === Source file: plugins/agent-tools.md ===
\section{Plugin agent tools}


OpenClaw plugins can register \textbf{agent tools} (JSON‑schema functions) that are exposed
to the LLM during agent runs. Tools can be \textbf{required} (always available) or
\textbf{optional} (opt‑in).

Agent tools are configured under \texttt{tools} in the main config, or per‑agent under
\texttt{agents.list[].tools}. The allowlist/denylist policy controls which tools the agent
can call.

\subsection{Basic tool}


\begin{lstlisting}[language=typescript]
import { Type } from "@sinclair/typebox";

export default function (api) {
  api.registerTool({
    name: "my_tool",
    description: "Do a thing",
    parameters: Type.Object({
      input: Type.String(),
    }),
    async execute(_id, params) {
      return { content: [{ type: "text", text: params.input }] };
    },
  });
}
\end{lstlisting}


\subsection{Optional tool (opt‑in)}


Optional tools are \textbf{never} auto‑enabled. Users must add them to an agent
allowlist.

\begin{lstlisting}[language=typescript]
export default function (api) {
  api.registerTool(
    {
      name: "workflow_tool",
      description: "Run a local workflow",
      parameters: {
        type: "object",
        properties: {
          pipeline: { type: "string" },
        },
        required: ["pipeline"],
      },
      async execute(_id, params) {
        return { content: [{ type: "text", text: params.pipeline }] };
      },
    },
    { optional: true },
  );
}
\end{lstlisting}


Enable optional tools in \texttt{agents.list[].tools.allow} (or global \texttt{tools.allow}):

\begin{lstlisting}[language=json]
{
  agents: {
    list: [
      {
        id: "main",
        tools: {
          allow: [
            "workflow_tool", // specific tool name
            "workflow", // plugin id (enables all tools from that plugin)
            "group:plugins", // all plugin tools
          ],
        },
      },
    ],
  },
}
\end{lstlisting}


Other config knobs that affect tool availability:

\begin{itemize}
  \item Allowlists that only name plugin tools are treated as plugin opt-ins; core tools remain
\end{itemize}

enabled unless you also include core tools or groups in the allowlist.
\begin{itemize}
  \item \texttt{tools.profile} / \texttt{agents.list[].tools.profile} (base allowlist)
  \item \texttt{tools.byProvider} / \texttt{agents.list[].tools.byProvider} (provider‑specific allow/deny)
  \item \texttt{tools.sandbox.tools.*} (sandbox tool policy when sandboxed)
\end{itemize}


\subsection{Rules + tips}


\begin{itemize}
  \item Tool names must \textbf{not} clash with core tool names; conflicting tools are skipped.
  \item Plugin ids used in allowlists must not clash with core tool names.
  \item Prefer \texttt{optional: true} for tools that trigger side effects or require extra
\end{itemize}

binaries/credentials.


\clearpage

% === Source file: plugins/community.md ===
\section{Community plugins}


This page tracks high-quality \textbf{community-maintained plugins} for OpenClaw.

We accept PRs that add community plugins here when they meet the quality bar.

\subsection{Required for listing}


\begin{itemize}
  \item Plugin package is published on npmjs (installable via \texttt{openclaw plugins install }).
  \item Source code is hosted on GitHub (public repository).
  \item Repository includes setup/use docs and an issue tracker.
  \item Plugin has a clear maintenance signal (active maintainer, recent updates, or responsive issue handling).
\end{itemize}


\subsection{How to submit}


Open a PR that adds your plugin to this page with:

\begin{itemize}
  \item Plugin name
  \item npm package name
  \item GitHub repository URL
  \item One-line description
  \item Install command
\end{itemize}


\subsection{Review bar}


We prefer plugins that are useful, documented, and safe to operate.
Low-effort wrappers, unclear ownership, or unmaintained packages may be declined.

\subsection{Candidate format}


Use this format when adding entries:

\begin{itemize}
  \item \textbf{Plugin Name} — short description
\end{itemize}

npm: \texttt{@scope/package}
repo: \texttt{https://github.com/org/repo}
install: \texttt{openclaw plugins install @scope/package}

\subsection{Listed plugins}


\begin{itemize}
  \item \textbf{WeChat} — Connect OpenClaw to WeChat personal accounts via WeChatPadPro (iPad protocol). Supports text, image, and file exchange with keyword-triggered conversations.
\end{itemize}

npm: \texttt{@icesword760/openclaw-wechat}
repo: \texttt{https://github.com/icesword0760/openclaw-wechat}
install: \texttt{openclaw plugins install @icesword760/openclaw-wechat}


\clearpage

% === Source file: plugins/manifest.md ===
\section{Plugin manifest (openclaw.plugin.json)}


Every plugin \textbf{must} ship a \texttt{openclaw.plugin.json} file in the \textbf{plugin root}.
OpenClaw uses this manifest to validate configuration **without executing plugin
code**. Missing or invalid manifests are treated as plugin errors and block
config validation.

See the full plugin system guide: \href{/tools/plugin}{Plugins}.

\subsection{Required fields}


\begin{lstlisting}[language=json]
{
  "id": "voice-call",
  "configSchema": {
    "type": "object",
    "additionalProperties": false,
    "properties": {}
  }
}
\end{lstlisting}


Required keys:

\begin{itemize}
  \item \texttt{id} (string): canonical plugin id.
  \item \texttt{configSchema} (object): JSON Schema for plugin config (inline).
\end{itemize}


Optional keys:

\begin{itemize}
  \item \texttt{kind} (string): plugin kind (example: \texttt{"memory"}).
  \item \texttt{channels} (array): channel ids registered by this plugin (example: \texttt{["matrix"]}).
  \item \texttt{providers} (array): provider ids registered by this plugin.
  \item \texttt{skills} (array): skill directories to load (relative to the plugin root).
  \item \texttt{name} (string): display name for the plugin.
  \item \texttt{description} (string): short plugin summary.
  \item \texttt{uiHints} (object): config field labels/placeholders/sensitive flags for UI rendering.
  \item \texttt{version} (string): plugin version (informational).
\end{itemize}


\subsection{JSON Schema requirements}


\begin{itemize}
  \item \textbf{Every plugin must ship a JSON Schema}, even if it accepts no config.
  \item An empty schema is acceptable (for example, \texttt{\{ "type": "object", "additionalProperties": false \}}).
  \item Schemas are validated at config read/write time, not at runtime.
\end{itemize}


\subsection{Validation behavior}


\begin{itemize}
  \item Unknown \texttt{channels.*} keys are \textbf{errors}, unless the channel id is declared by
\end{itemize}

a plugin manifest.
\begin{itemize}
  \item \texttt{plugins.entries.}, \texttt{plugins.allow}, \texttt{plugins.deny}, and \texttt{plugins.slots.*}
\end{itemize}

must reference \textbf{discoverable} plugin ids. Unknown ids are \textbf{errors}.
\begin{itemize}
  \item If a plugin is installed but has a broken or missing manifest or schema,
\end{itemize}

validation fails and Doctor reports the plugin error.
\begin{itemize}
  \item If plugin config exists but the plugin is \textbf{disabled}, the config is kept and
\end{itemize}

a \textbf{warning} is surfaced in Doctor + logs.

\subsection{Notes}


\begin{itemize}
  \item The manifest is \textbf{required for all plugins}, including local filesystem loads.
  \item Runtime still loads the plugin module separately; the manifest is only for
\end{itemize}

discovery + validation.
\begin{itemize}
  \item If your plugin depends on native modules, document the build steps and any
\end{itemize}

package-manager allowlist requirements (for example, pnpm \texttt{allow-build-scripts}
\begin{itemize}
  \item \texttt{pnpm rebuild }).
\end{itemize}



\clearpage

% === Source file: plugins/voice-call.md ===
\section{Voice Call (plugin)}


Voice calls for OpenClaw via a plugin. Supports outbound notifications and
multi-turn conversations with inbound policies.

Current providers:

\begin{itemize}
  \item \texttt{twilio} (Programmable Voice + Media Streams)
  \item \texttt{telnyx} (Call Control v2)
  \item \texttt{plivo} (Voice API + XML transfer + GetInput speech)
  \item \texttt{mock} (dev/no network)
\end{itemize}


Quick mental model:

\begin{itemize}
  \item Install plugin
  \item Restart Gateway
  \item Configure under \texttt{plugins.entries.voice-call.config}
  \item Use \texttt{openclaw voicecall ...} or the \texttt{voice\_call} tool
\end{itemize}


\subsection{Where it runs (local vs remote)}


The Voice Call plugin runs \textbf{inside the Gateway process}.

If you use a remote Gateway, install/configure the plugin on the \textbf{machine running the Gateway}, then restart the Gateway to load it.

\subsection{Install}


\subsubsection{Option A: install from npm (recommended)}


\begin{lstlisting}[language=bash]
openclaw plugins install @openclaw/voice-call
\end{lstlisting}


Restart the Gateway afterwards.

\subsubsection{Option B: install from a local folder (dev, no copying)}


\begin{lstlisting}[language=bash]
openclaw plugins install ./extensions/voice-call
cd ./extensions/voice-call && pnpm install
\end{lstlisting}


Restart the Gateway afterwards.

\subsection{Config}


Set config under \texttt{plugins.entries.voice-call.config}:

\begin{lstlisting}[language=json]
{
  plugins: {
    entries: {
      "voice-call": {
        enabled: true,
        config: {
          provider: "twilio", // or "telnyx" | "plivo" | "mock"
          fromNumber: "+15550001234",
          toNumber: "+15550005678",

          twilio: {
            accountSid: "ACxxxxxxxx",
            authToken: "...",
          },

          telnyx: {
            apiKey: "...",
            connectionId: "...",
            // Telnyx webhook public key from the Telnyx Mission Control Portal
            // (Base64 string; can also be set via TELNYX_PUBLIC_KEY).
            publicKey: "...",
          },

          plivo: {
            authId: "MAxxxxxxxxxxxxxxxxxxxx",
            authToken: "...",
          },

          // Webhook server
          serve: {
            port: 3334,
            path: "/voice/webhook",
          },

          // Webhook security (recommended for tunnels/proxies)
          webhookSecurity: {
            allowedHosts: ["voice.example.com"],
            trustedProxyIPs: ["100.64.0.1"],
          },

          // Public exposure (pick one)
          // publicUrl: "https://example.ngrok.app/voice/webhook",
          // tunnel: { provider: "ngrok" },
          // tailscale: { mode: "funnel", path: "/voice/webhook" }

          outbound: {
            defaultMode: "notify", // notify | conversation
          },

          streaming: {
            enabled: true,
            streamPath: "/voice/stream",
            preStartTimeoutMs: 5000,
            maxPendingConnections: 32,
            maxPendingConnectionsPerIp: 4,
            maxConnections: 128,
          },
        },
      },
    },
  },
}
\end{lstlisting}


Notes:

\begin{itemize}
  \item Twilio/Telnyx require a \textbf{publicly reachable} webhook URL.
  \item Plivo requires a \textbf{publicly reachable} webhook URL.
  \item \texttt{mock} is a local dev provider (no network calls).
  \item Telnyx requires \texttt{telnyx.publicKey} (or \texttt{TELNYX\_PUBLIC\_KEY}) unless \texttt{skipSignatureVerification} is true.
  \item \texttt{skipSignatureVerification} is for local testing only.
  \item If you use ngrok free tier, set \texttt{publicUrl} to the exact ngrok URL; signature verification is always enforced.
  \item \texttt{tunnel.allowNgrokFreeTierLoopbackBypass: true} allows Twilio webhooks with invalid signatures \textbf{only} when \texttt{tunnel.provider="ngrok"} and \texttt{serve.bind} is loopback (ngrok local agent). Use for local dev only.
  \item Ngrok free tier URLs can change or add interstitial behavior; if \texttt{publicUrl} drifts, Twilio signatures will fail. For production, prefer a stable domain or Tailscale funnel.
  \item Streaming security defaults:
  \item \texttt{streaming.preStartTimeoutMs} closes sockets that never send a valid \texttt{start} frame.
  \item \texttt{streaming.maxPendingConnections} caps total unauthenticated pre-start sockets.
  \item \texttt{streaming.maxPendingConnectionsPerIp} caps unauthenticated pre-start sockets per source IP.
  \item \texttt{streaming.maxConnections} caps total open media stream sockets (pending + active).
\end{itemize}


\subsection{Stale call reaper}


Use \texttt{staleCallReaperSeconds} to end calls that never receive a terminal webhook
(for example, notify-mode calls that never complete). The default is \texttt{0}
(disabled).

Recommended ranges:

\begin{itemize}
  \item \textbf{Production:} \texttt{120}–\texttt{300} seconds for notify-style flows.
  \item Keep this value \textbf{higher than \texttt{maxDurationSeconds}} so normal calls can
\end{itemize}

finish. A good starting point is \texttt{maxDurationSeconds + 30–60} seconds.

Example:

\begin{lstlisting}[language=json]
{
  plugins: {
    entries: {
      "voice-call": {
        config: {
          maxDurationSeconds: 300,
          staleCallReaperSeconds: 360,
        },
      },
    },
  },
}
\end{lstlisting}


\subsection{Webhook Security}


When a proxy or tunnel sits in front of the Gateway, the plugin reconstructs the
public URL for signature verification. These options control which forwarded
headers are trusted.

\texttt{webhookSecurity.allowedHosts} allowlists hosts from forwarding headers.

\texttt{webhookSecurity.trustForwardingHeaders} trusts forwarded headers without an allowlist.

\texttt{webhookSecurity.trustedProxyIPs} only trusts forwarded headers when the request
remote IP matches the list.

Webhook replay protection is enabled for Twilio and Plivo. Replayed valid webhook
requests are acknowledged but skipped for side effects.

Twilio conversation turns include a per-turn token in `` callbacks, so
stale/replayed speech callbacks cannot satisfy a newer pending transcript turn.

Example with a stable public host:

\begin{lstlisting}[language=json]
{
  plugins: {
    entries: {
      "voice-call": {
        config: {
          publicUrl: "https://voice.example.com/voice/webhook",
          webhookSecurity: {
            allowedHosts: ["voice.example.com"],
          },
        },
      },
    },
  },
}
\end{lstlisting}


\subsection{TTS for calls}


Voice Call uses the core \texttt{messages.tts} configuration (OpenAI or ElevenLabs) for
streaming speech on calls. You can override it under the plugin config with the
\textbf{same shape} — it deep‑merges with \texttt{messages.tts}.

\begin{lstlisting}[language=json]
{
  tts: {
    provider: "elevenlabs",
    elevenlabs: {
      voiceId: "pMsXgVXv3BLzUgSXRplE",
      modelId: "eleven_multilingual_v2",
    },
  },
}
\end{lstlisting}


Notes:

\begin{itemize}
  \item \textbf{Edge TTS is ignored for voice calls} (telephony audio needs PCM; Edge output is unreliable).
  \item Core TTS is used when Twilio media streaming is enabled; otherwise calls fall back to provider native voices.
\end{itemize}


\subsubsection{More examples}


Use core TTS only (no override):

\begin{lstlisting}[language=json]
{
  messages: {
    tts: {
      provider: "openai",
      openai: { voice: "alloy" },
    },
  },
}
\end{lstlisting}


Override to ElevenLabs just for calls (keep core default elsewhere):

\begin{lstlisting}[language=json]
{
  plugins: {
    entries: {
      "voice-call": {
        config: {
          tts: {
            provider: "elevenlabs",
            elevenlabs: {
              apiKey: "elevenlabs_key",
              voiceId: "pMsXgVXv3BLzUgSXRplE",
              modelId: "eleven_multilingual_v2",
            },
          },
        },
      },
    },
  },
}
\end{lstlisting}


Override only the OpenAI model for calls (deep‑merge example):

\begin{lstlisting}[language=json]
{
  plugins: {
    entries: {
      "voice-call": {
        config: {
          tts: {
            openai: {
              model: "gpt-4o-mini-tts",
              voice: "marin",
            },
          },
        },
      },
    },
  },
}
\end{lstlisting}


\subsection{Inbound calls}


Inbound policy defaults to \texttt{disabled}. To enable inbound calls, set:

\begin{lstlisting}[language=json]
{
  inboundPolicy: "allowlist",
  allowFrom: ["+15550001234"],
  inboundGreeting: "Hello! How can I help?",
}
\end{lstlisting}


Auto-responses use the agent system. Tune with:

\begin{itemize}
  \item \texttt{responseModel}
  \item \texttt{responseSystemPrompt}
  \item \texttt{responseTimeoutMs}
\end{itemize}


\subsection{CLI}


\begin{lstlisting}[language=bash]
openclaw voicecall call --to "+15555550123" --message "Hello from OpenClaw"
openclaw voicecall continue --call-id <id> --message "Any questions?"
openclaw voicecall speak --call-id <id> --message "One moment"
openclaw voicecall end --call-id <id>
openclaw voicecall status --call-id <id>
openclaw voicecall tail
openclaw voicecall expose --mode funnel
\end{lstlisting}


\subsection{Agent tool}


Tool name: \texttt{voice\_call}

Actions:

\begin{itemize}
  \item \texttt{initiate\_call} (message, to?, mode?)
  \item \texttt{continue\_call} (callId, message)
  \item \texttt{speak\_to\_user} (callId, message)
  \item \texttt{end\_call} (callId)
  \item \texttt{get\_status} (callId)
\end{itemize}


This repo ships a matching skill doc at \texttt{skills/voice-call/SKILL.md}.

\subsection{Gateway RPC}


\begin{itemize}
  \item \texttt{voicecall.initiate} (\texttt{to?}, \texttt{message}, \texttt{mode?})
  \item \texttt{voicecall.continue} (\texttt{callId}, \texttt{message})
  \item \texttt{voicecall.speak} (\texttt{callId}, \texttt{message})
  \item \texttt{voicecall.end} (\texttt{callId})
  \item \texttt{voicecall.status} (\texttt{callId})
\end{itemize}



\clearpage

% === Source file: plugins/zalouser.md ===
\section{Zalo Personal (plugin)}


Zalo Personal support for OpenClaw via a plugin, using native \texttt{zca-js} to automate a normal Zalo user account.

\begin{quote}
\textbf{Warning:} Unofficial automation may lead to account suspension/ban. Use at your own risk.
\end{quote}


\subsection{Naming}


Channel id is \texttt{zalouser} to make it explicit this automates a \textbf{personal Zalo user account} (unofficial). We keep \texttt{zalo} reserved for a potential future official Zalo API integration.

\subsection{Where it runs}


This plugin runs \textbf{inside the Gateway process}.

If you use a remote Gateway, install/configure it on the \textbf{machine running the Gateway}, then restart the Gateway.

No external \texttt{zca}/\texttt{openzca} CLI binary is required.

\subsection{Install}


\subsubsection{Option A: install from npm}


\begin{lstlisting}[language=bash]
openclaw plugins install @openclaw/zalouser
\end{lstlisting}


Restart the Gateway afterwards.

\subsubsection{Option B: install from a local folder (dev)}


\begin{lstlisting}[language=bash]
openclaw plugins install ./extensions/zalouser
cd ./extensions/zalouser && pnpm install
\end{lstlisting}


Restart the Gateway afterwards.

\subsection{Config}


Channel config lives under \texttt{channels.zalouser} (not \texttt{plugins.entries.*}):

\begin{lstlisting}[language=json]
{
  channels: {
    zalouser: {
      enabled: true,
      dmPolicy: "pairing",
    },
  },
}
\end{lstlisting}


\subsection{CLI}


\begin{lstlisting}[language=bash]
openclaw channels login --channel zalouser
openclaw channels logout --channel zalouser
openclaw channels status --probe
openclaw message send --channel zalouser --target <threadId> --message "Hello from OpenClaw"
openclaw directory peers list --channel zalouser --query "name"
\end{lstlisting}


\subsection{Agent tool}


Tool name: \texttt{zalouser}

Actions: \texttt{send}, \texttt{image}, \texttt{link}, \texttt{friends}, \texttt{groups}, \texttt{me}, \texttt{status}

Channel message actions also support \texttt{react} for message reactions.


\clearpage
